\section{Introduction}
\label{sec:introduction}

\subsection{Background and Motivation}
Automated credit scoring supports lending decisions in which the two error types are
economically asymmetric: approving an applicant who turns out to default (a missed
\texttt{bad} case) destroys principal, whereas rejecting a creditworthy applicant (a
false \texttt{bad} alarm) merely forgoes interest income. The Statlog German Credit
data makes this asymmetry explicit through an official 5:1 misclassification-cost
matrix, in which missing a \texttt{bad} applicant costs five times as much as
rejecting a \texttt{good} one~\cite{uci_credit_g_1994}. Class imbalance amplifies the
problem: \texttt{bad} applicants form only 30\% of \texttt{credit-g} and about 6\% of
the \texttt{sick} thyroid subset~\cite{uci_thyroid_1986}, so accuracy is dominated by
the majority class --- a trivial rule that always predicts \texttt{good} already
achieves 0.70 accuracy on \texttt{credit-g} while never identifying a single default.
Practitioners can draw on three families of remedies --- synthetic minority
over-sampling~\cite{chawla2002smote}, cost-sensitive training, and decision-threshold
adjustment --- yet their relative merits change with the model family and the degree
of imbalance, and are rarely compared under one leakage-safe protocol with an
application-grounded cost measure.

\subsection{Research Significance}
This project delivers such a comparison on two UCI datasets~\cite{dua2019uci} that
differ deliberately in scale, imbalance ratio, missingness and domain. For risk
management, it shows which imbalance intervention actually minimises the official
misclassification cost rather than a surrogate metric. For methodology, it
demonstrates a reproducible experimental discipline: preprocessing specifications
frozen before the main matrix, locked data splits, SHA-256 provenance chained from
raw data to every result file, and all imputation, encoding, scaling and resampling
fitted strictly inside the cross-validation folds, so that any observed gain is
attributable to the strategies rather than to information leakage. For
interpretability, association rules mined from the same data are triangulated
against tree-based model explanations. The full pipeline is implemented with
scikit-learn~\cite{pedregosa2011sklearn}, imbalanced-learn~\cite{lemaitre2017imblearn}
and mlxtend~\cite{raschka2018mlxtend}, with Weka~\cite{frank2016weka} as the course
reference tool.

\subsection{Research Questions and Contributions}
The study is organised around four research questions:
\begin{itemize}
  \item \textbf{RQ1.} How does the degree of class imbalance moderate the gains of
  imbalance-handling strategies? The two datasets provide a 70:30 and a 94:6
  endpoint for this question.
  \item \textbf{RQ2.} Under what conditions does accuracy mislead model selection on
  imbalanced data, and which metrics (positive-class F1, ROC-AUC, official
  misclassification cost) track decision utility instead?
  \item \textbf{RQ3.} Under the official 5:1 cost matrix, does the best
  model--strategy pairing change relative to the rankings obtained by accuracy or F1?
  \item \textbf{RQ4.} How does the encoding of nominal attributes (one-hot vs.\
  ordinal) interact with the model family (linear/kernel vs.\ tree-based) on
  mixed-attribute data?
\end{itemize}

\noindent The main contributions are: (i) a paired two-dataset design contrasting
five dimensions at once --- scale (1{,}000 vs.\ 3{,}772 records), imbalance (70:30
vs.\ 94:6), missingness (0 vs.\ about 16\%), attribute composition and domain ---
which directly probes RQ1; (ii) a full experiment matrix of four classifiers
(decision tree, random forest, logistic regression, RBF-SVM) $\times$ four
strategies (plain, SMOTE, cost-sensitive, threshold adjustment) built on
leakage-safe pipelines with locked splits; (iii) the official cost matrix treated
as a first-class evaluation axis alongside F1 and ROC-AUC (RQ2, RQ3); (iv) a
controlled encoding ablation with a supplementary attribution control that
separates encoding from scaling effects (RQ4); and (v) an interpretability
triangulation of Apriori rules, decision-tree paths and random-forest feature
importances (Section~\ref{sec:apriori}). Statistical comparisons use fold-level
tests with the Nadeau--Bengio
correction~\cite{nadeau2003inference,demsar2006comparison}. Experimental results
are reported in Section~\ref{sec:results}.
